\section{Admissible cuts and derived equivalences}

This section contains the main result of the paper. We show how, starting from a finite poset $P$, one can recognize directly from its order structure decompositions to which Theorem \ref{ladkani} can be applied, and hence construct new posets universally derived equivalent to $P$. More precisely, given a partition $P=A \sqcup B$, we characterize when the data required by that theorem are already encoded in the order structure of $P$. This leads to the intrinsic notion of an admissible cut.

Once such an admissible cut is found, the required data are recovered from $P$ itself and determine an associated poset $P^-$ universally derived equivalent to $P$. We show that this transformation is reversible once the partition $A \sqcup B$ is retained, and we illustrate the construction with the Tamari lattice $T_3$. We then derive an application to posets of height at most one, study the singleton specialization and some constraints on the cardinalities of the sets $M(a)$ and finally record an invariance property under isomorphisms. 

Throughout this section, $P$ denotes a finite poset and $P=A \sqcup B$ a partition of its underlying set. Whenever needed, $A$ and $B$ are endowed with the partial orders induced from $P$, denoted by $\leq_A$ and $\leq_B$, respectively.

\subsection{The general construction}

Recall that a subset $A \subseteq P$ is called an \textbf{order ideal} if, whenever $a \in A$ and $ x \leq_P a $, then $ x \in A$. Dually, a subset $B \subseteq P$ is called an \textbf{order filter} if, whenever $b \in B $ and $ b \leq_P x$, then $x \in B$. From now on, we simply write ideals and filters.

With the above notation, for a subset $S \subseteq P$, define
\[U_B(S) \coloneq \{b \in B \colon \text{ there exists } s \in S \text{ such that } s \leq_Pb\}.\] When $S=\{x\}$ is a singleton, we write $U_B(x)$ instead of $U_B(\{x\})$. If $S \subseteq B$, then $U_B(S)$ is the upward closure of $S$ in $B$.

For a subset $C \subseteq B$, let $\operatorname{Min}_B(C)$ and $\operatorname{Max}_B(C)$ denote the sets of minimal and maximal elements of $C$, respectively, with respect to $\leq_B$. For every $a \in A$, set
\[M(a) \coloneq \operatorname{Min}_B(U_B(a)).\]

For $ b\in B$, we write 
\[[b, \bullet]_B \coloneq \{b' \in B \colon b \leq_Bb'\}, \,\,\,\,\, [\bullet,b]_B \coloneq \{b'\in B \colon b' \leq_Bb\}.\]

\begin{definition} %[Admissible cut] 
\label{def admissible cuts}
   Let $P$ be a finite poset and let $P=A \sqcup B$ be a partition of its underlying set. 
   We say that $P=A \sqcup B$ is an \textbf{admissible cut} if the following conditions are satisfied:
\begin{enumerate}
    \item [(a)] Ideal-filter condition: $A$ is an ideal of $P$ and $B$ is a filter of $P$;
    \item [(b)]  Separation condition: for every $ a \in A$ and every pair of distinct elements $m, m' \in M(a)$, \[[m, \bullet]_B \cap [m', \bullet]_B = \emptyset \,\,\,\, \text{ and } \,\,\,\, [\bullet, m]_B \cap [\bullet, m']_B = \emptyset.\]
    \item [(c)] Compatibility condition: for every $a \leq_A a'$ and every $m \in M(a)$, there exists $m'\in~M(a')$ such that $m \leq_B m'$.
\end{enumerate}
\end{definition}

\begin{theorem} %[Recognition theorem for admissible cuts]
\label{generalizzazione teomio}
    Let $P$ be a finite poset, and let $P=A \sqcup B$ be a partition of its underlying set.

     The following are equivalent: 
     \begin{enumerate}
    \item There exists a family $\mathcal{B}=(B_a)_{a \in A}$ of subsets of $B$ satisfying the hypotheses of Ladkani's construction such that $P=(A \sqcup B)^+_{\mathcal{B}}$.
    
    \item $P=A \sqcup B$ is an admissible cut.
    \end{enumerate}
   Moreover, whenever these equivalent conditions hold, the subsets $B_a$ in $(i)$ are necessarily given by 
   \[B_a=M(a) \,\, \text{ for every } \,\, a \in A.\]

\end{theorem}

\begin{proof}
Let $P$ be a finite poset.
\item[\boxed{(i)\Rightarrow(ii)}:]
Assume that a family $(B_a)_{a \in A}$ as in $(i)$ exists. By assumption, for every $a \in A$ and $b \in B$, \begin{equation} \label{defP+}
 a \leq_P b \iff \text{ there exists }\, y \in B_a \text{ such that } y \leq_B b.
 \end{equation}
 \begin{itemize}
 \item \textbf{\underline{Step 1: Ideal-filter condition.}} 
 
 In the poset $(A \sqcup B, \leq_+)$ the only mixed  relations are of the form \[a \leq_+b, \,\,\,\,\,\,\,\, a \in A,\,\, b \in B.\]
 
  If $x \leq_P a$ with $a \in A$, then $x$ cannot belong to $B$. Hence $x \in A$, and therefore $A$ is an ideal of $P$. Dually, if $b \leq_Px$ with $b \in B$, then $x$ cannot belong to $A$. Hence $x \in B$, and therefore $B$ is a filter of $P$.
 
\item \textbf{\underline{Step 2: $B_a=M(a)$ for every $a \in A$.}}

Fix $a \in A$. By (\ref{defP+}), we have \[U_B(a)=\{b \in B \colon a \leq_P b\}=\{b \in B \colon \text{ there exists } \, y \in B_a \text{ such that } y \leq_B b\}=U_B(B_a).\]
We first prove that $B_a \subseteq M(a)$. Let $y \in B_a$. Since $y \leq_By$, we have $y \in U_B(B_a)=U_B(a)$. Suppose that $z \in U_B(a)$ and $z \leq_B y$. Since $z \in U_B(a)=U_B(B_a)$, there exists $y' \in B_a$ such that $y' \leq_B z \leq_B y$. If $y' \neq y$, then $y \in [y', \bullet]_B \cap [y, \bullet]_B$, contradicting condition (i) of the construction in Section \ref{Ladkani's construction}, applied to the distinct elements $y,y' \in B_a$. Hence $y'=y$, and from $y \leq_B z \leq_B y$ we get $z=y$. So $y$ is minimal in $U_B(a)$, i.e. $y \in M(a)$.

Conversely, let $z \in M(a)$. Since $z \in U_B(a)= U_B(B_a)$, there exists $y \in B_a$ such that $y \leq_Bz$. By the inclusion already proved, $y \in M(a)$. Since $z$ is minimal in $U_B(a)$, while $y \in U_B(a)$ and $y \leq_B z$, we obtain $y=z$. Hence $z \in B_a$.

Therefore, $B_a=M(a)$.

\item \textbf{\underline{Step 3: Separation condition.}} Since $B_a=M(a)$ for every $a \in A$, condition (b) is exactly condition (i) of the construction in Section \ref{Ladkani's construction} applied to the family $(B_a)_{a \in A}$.

\item \textbf{\underline{Step 4: Compatibility condition.}} Let $a \leq a'$ in $A$. By condition (ii) of the construction in Section \ref{Ladkani's construction}, there exists a bijection $\varphi_{a,a'} \colon B_a \to B_{a'}$ such that $y \leq_B \varphi_{a,a'}(y)$ for every $y \in B_a$. Since $B_a=M(a)$ and $B_{a'}=M(a')$, for every $m \in M(a)$ the element $\varphi_{a,a'}(m)$ belongs to $M(a')$ and satisfies $m \leq_B \varphi_{a,a'}(m)$. Thus condition (c) holds.
\end{itemize}

Therefore, the partition $P=A \sqcup B$ satisfies the conditions $(a)$, $(b)$ and $(c)$, and is an admissible cut.
 
 \item[\boxed{(ii)\Rightarrow(i)}:] Assume that $P=A \sqcup B$ is an admissible cut. Define, for each $a  \in A$, \[B_a \coloneq M(a).\]
 We show that $(B_a)_{a \in A}$ satisfies the hypotheses of Section \ref{Ladkani's construction} and that the resulting positive poset coincides with $P$. 
\begin{itemize}
 \item \textbf{\underline{Step 1: Verification of condition (i).}} By definition $B_a=M(a)$. For every $a \in A$ and every $m \neq m' \in B_a=M(a)$,
 \[[m, \bullet]_B \cap [m', \bullet]_B = \emptyset, \,\,\,\,\, [\bullet,m]_B \cap [\bullet,m']_B = \emptyset.\]
 \item \textbf{\underline{Step 2: Verification of condition (ii).}} Let $a \leq a'$ in $A$. 
 
 We construct $\varphi_{a,a'} \colon M(a) \to M(a')$ with $m \leq_B \varphi_{a,a'}(m)$ for every $m \in M(a)$, and show it is a bijection.

 By condition (c), for each $m \in M(a)$ there exists at least one $m' \in M(a')$ with $m \leq_Bm'$.

 \textit{Uniqueness of $m'$}: suppose $m \leq_B m'_1$ and $m \leq_B m_2'$ with $m_1',m_2' \in M(a')$. Then $m \in [\bullet,m_1']_B \cap [\bullet,m_2']_B$. 
 
 By condition (b) applied to $a'$, this intersection is empty whenever $m_1' \neq m_2'$. Hence $m_1' = m_2'$.

 Then we set $\varphi_{a,a'}(m) \coloneq m'$, where $m'$ is the unique element in $M(a')$ with $m \leq_Bm'$.
 By construction, $m \leq_B \varphi_{a,a'}(m)$ for every $m \in M(a)$.
 
 \textit{Injectivity:} Suppose $\varphi_{a,a'}(m_1)=\varphi_{a,a'}(m_2)=m'$. Then $m' \in [m_1, \bullet]_B \cap [m_2, \bullet]_B$. By condition (b) applied to $a$, this intersection is empty whenever $m_1 \neq m_2$. Hence $m_1=m_2$.

 \textit{Surjectivity:} Let $m' \in M(a')$. From $a \leq_P a' \leq_Pm'$, we have $m' \in U_B(a)$. Since $P$ is finite, every element of a nonempty subset lies above at least one minimal element of that subset. Applied to $U_B(a)$, this gives an element $m \in M(a)$ such that $m \leq_Bm'$. By uniqueness, $\varphi_{a,a'}(m)=m'$. 

 Therefore $\varphi_{a,a'}$ is a bijection satisfying \[m \leq_B \varphi_{a,a'}(m)\] for every $m \in M(a)$. Hence the family $(B_a)_{a \in A}$, with $B_a=M(a)$, satisfies condition $(ii)$ of Section \ref{Ladkani's construction}.

 \item \textbf{\underline{Step 3: Reconstruction of the mixed relations.}}
 Let $P^+=(A \sqcup B, \leq_+)$ be the positive poset associated with the family $B_a=M(a)$. The orders inside $A$ and $B$ coincide in $P$ and $P^+$ by construction, so it suffices to compare the mixed relations.

 Let $a \in A$ and $b \in B$. By definition of $\leq_+$, 
 \begin{equation} \label{defP+M(a)} 
 a \leq_+b \iff \exists m \in M(a) \text{ with } m \leq_Bb.
 \end{equation}
 Since $U_B(a)$ is upward closed in $B$ and $B$ is finite, every element of $U_B(a)$ lies above a minimal element of $U_B(a)$. Since \[M(a) = \operatorname{Min}_B(U_B(a)),\] it follows that \[U_B(a)=U_B(M(a)).\]

 Combining (\ref{defP+M(a)}) with the equality $U_B(a)=U_B(M(a))$, we obtain, for every $a \in A$ and $b \in B$, \[a \leq_+b \iff b \in U_B(M(a)) \iff b \in U_B(a) \iff a \leq_Pb.\]

 Finally, there are no mixed relations of the form $b \leq a$ (with $b \in B, a \in A$) in $P$ (because $A$ is an ideal by (a)) nor in $P^+$ (by construction). Hence $P$ and $P^+$ have the same partial order, i.e., $P^+=P$.

 \end{itemize}
\end{proof}
For a fixed admissible cut $P=A \sqcup B$, let $\mathcal{M}=(M(a))_{a \in A}$. We denote the positive and negative posets associated with this family by \[P^+ \coloneq (A \sqcup B)^+_{\mathcal{M}}, \,\,\,\,\,\,\, P^- \coloneq (A \sqcup B)^-_{\mathcal{M}}.\] By Theorem \ref{generalizzazione teomio}, $P^+=P$.

Since $P$ is finite, Theorem \ref{generalizzazione teomio} gives a finite procedure for detecting admissible cuts. For a given partition $P=A \sqcup B$, one computes the sets $M(a)$ and checks conditions (a)--(c). By running this test over the finitely many partitions of the underlying set of $P$, one can determine all admissible cuts of $P$ and hence all associated negative posets $P^-$.

 The following corollary describes $P^-$ explicitly.

\begin{corollary} %[The associated universally derived equivalent poset] 
\label{The associated universally derived equivalent poset}
Let $P$ be a finite poset and let $P=A \sqcup B$ be an admissible cut. Then the mixed relations of the associated negative poset $P^-$ are given by
\[b \leq_-a \iff \text{ there exists } m \in M(a) \text{ such that } b \leq_B m.\]
Moreover, $P^-$ is universally derived equivalent to $P$.
\end{corollary}

\begin{proof}
The description of the mixed relations follows directly from the definition of $P^-$.
By Theorem \ref{generalizzazione teomio}, $P^+=P$. By Lemma \ref{lemma part ord}, $P^-$ is a poset, and by Theorem \ref{ladkani}, $P^+$ and $P^-$ are universally derived equivalent. Hence $P$ and $P^-$ are universally derived equivalent.
\end{proof}

The next result shows that the construction is reversible once the partition $A \sqcup B$ is retained. More precisely, the data $M(a)$, and hence the original mixed relations of $P$, can be recovered from the associated negative poset $P^-$.

\begin{proposition}%[Dual reversibility of the construction/reconstruction from the negative poset] 
\label{Reversibility of the construction}

Let $P$ be a finite poset, let $P=A \sqcup B$ be an admissible cut, and let $P^-$ be the associated negative poset. Then, in $P^-$, $B$ is an ideal and $A$ is a filter. For every $a \in A$, set \[L_B^-(a) \coloneq \{b \in B \colon b \leq_-a\}.\]
Then
\[M(a)=\operatorname{Max}_B(L_B^-(a)).\]
Consequently, the original poset $P$ is uniquely determined by $P^-$ together with the partition $A \sqcup B$. More precisely, the orders induced on $A$ and $B$ are unchanged, and, for every $a \in A$ and $b \in B$,
\[a \leq_Pb \iff \text{ there exists } m \in \operatorname{Max}_B(L^-_B(a)) \text{ such that } m \leq_B b.\]
    
\end{proposition}

\begin{proof}
By the definition of the negative order, the only mixed relations in $P^-$ are of the form \[b \leq_-a, \,\,\,\, b \in B, \, a \in A.\]
Hence no element of $A$ lies below an element of $B$ in $P^-$. It follows that $B$ is an ideal of $P^-$. Dually, $A$ is a filter of $P^-$.

Now fix $a \in A$. By the definition of the negative order, \[L^-_B(a)= \bigcup_{m \in M(a)}[\bullet,m]_B.\]
    We first prove that every maximal element of $L^-_B(a)$ belongs to $M(a)$. Let $b\in L_B^-(a)$ be maximal. Then there exists $m \in M(a)$ such that $b \leq_Bm$. Since $m \in L_B^-(a)$, the maximality of $b$ implies $b=m$. Therefore
    \[\operatorname{Max}_B(L^-_B(a)) \subseteq M(a).\] Conversely, let $m \in M(a)$. Then $m \in L_B^-(a)$. Suppose that $m \leq_Bb$ for some $b \in L_B^-(a)$. By the definition of $L^-_B(a)$, there exists $m' \in M(a)$ such that $b\leq_B m'$. Hence $m \leq_Bb\leq_Bm'.$ Since $m,m' \in M(a)$, and $M(a)$ consists of the minimal elements of $U_B(a)$, the relation $m \leq_Bm'$ implies $m=m'$. It follows that $b=m$. Hence $m$ is maximal in $L^-_B(a)$, and therefore \[M(a) \subseteq \operatorname{Max}_B(L_B^-(a)).\]
    This proves \[M(a)=\operatorname{Max}_B(L^-_B(a)).\]
    The passage from $P$ to $P^-$ does not change the orders induced on $A$ and $B$. Moreover, since $A$ is an ideal of $P$, there are no mixed relations of the form $b \leq_Pa$, with $b \in B$ and $a \in A$.

    It remains to recover the mixed relations from $A$ to $B$. Since $B$ is finite and $U_B(a)$ is upward closed, it is the upward closure of its minimal elements. Hence
    \[a \leq_Pb \iff b \in U_B(a) \iff \text{ there exists } m \in M(a) \text{ such that } m\leq_Bb.\]
    Using the equality proved above, \[M(a)=\operatorname{Max}_B(L_B^-(a)),\] we obtain  
    \[a \leq_Pb \iff \text{ there exists } m \in \operatorname{Max}_B(L^-_B(a)) \text{ such that } m \leq_B b.\]
    Thus all the relations of $P$ can be reconstructed from $P^-$ and the partition $A \sqcup B$.
\end{proof}

Thus, once the partition $A \sqcup B$ is retained, passing from $P$ to $P^-$ does not lose any information. Notice that the roles of the two parts are reversed: $A$ is an ideal and $B$ is a filter in $P$, whereas $B$ is an ideal and $A$ is a filter in $P^-$. Correspondingly, the elements of $M(a)$, which are the minimal elements of $B$ lying above $a$ in $P$, are recovered as the maximal elements of $B$ lying below $a$ in $P^-$. In this sense, the reconstruction is dual to the original construction.

This duality can also be expressed by passing to opposite posets. Indeed, if $Q=(P^-)^{\operatorname{op}},$ then $A$ is an ideal and $B$ is a filter in $Q$, and for every $a \in A$, 
\[\operatorname{Min}_{B^{\operatorname{op}}}\{b \in B \colon a \leq_Qb\}= \operatorname{Max}_B(L^-_B(a))=M(a).\]
Thus, the data recovered from $P^-$ are precisely the order-dual counterparts of the data used in the original construction of Theorem \ref{generalizzazione teomio}.

\begin{example} \label{T3/D5}
    Consider the five-element Tamari lattice $T_3$ \cite{Geyer}, labelled as in the diagram on the left. Let 
    \[A=\{1,2,3,4\}, \,\,\, B=\{5\}.\]
    Since $5$ is the greatest element of $T_3$, the subset $A$ is an ideal and $B$ is a filter. Moreover, for every $a \in A$,
    \[U_B(a)=\{5\}, \,\,\, \text{ and hence }\,\,\, M(a)=\{5\}.\]
    The separation condition is vacuous, while the compatibility condition is immediate. Therefore $T_3=A \sqcup B$ is an admissible cut.
    % https://q.uiver.app/#q=WzAsMTAsWzIsMCwiNSBcXGJ1bGxldCJdLFswLDEsIjQgXFxidWxsZXQiXSxbNCwxLCIzIFxcYnVsbGV0Il0sWzAsMywiMiBcXGJ1bGxldCJdLFsyLDQsIjEgXFxidWxsZXQiXSxbOCwxLCI0IFxcYnVsbGV0Il0sWzEwLDEsIjJcXGJ1bGxldCJdLFsxMiwxLCIxIFxcYnVsbGV0Il0sWzE0LDEsIjMgXFxidWxsZXQiXSxbMTIsMywiNSBcXGJ1bGxldCJdLFsxLDBdLFsyLDBdLFszLDFdLFs0LDNdLFs0LDJdLFs2LDVdLFs3LDZdLFs3LDhdLFs5LDddXQ==
\[\begin{tikzcd}[sep=small]
	&& {5 \bullet} &&&&&&&&&&&& \\
	{4 \bullet} &&&& {3 \bullet} &&&& {4 \bullet} && {2\bullet} && {1 \bullet} && {3 \bullet} \\
	\\
	{2 \bullet} &&&&&&&&&&&& {5 \bullet} \\
	&& {1 \bullet}
	\arrow[from=2-1, to=1-3]
	\arrow[from=2-5, to=1-3]
	\arrow[from=2-11, to=2-9]
	\arrow[from=2-13, to=2-11]
	\arrow[from=2-13, to=2-15]
	\arrow[from=4-1, to=2-1]
	\arrow[from=4-13, to=2-13]
	\arrow[from=5-3, to=2-5]
	\arrow[from=5-3, to=4-1]
\end{tikzcd}\]
The poset displayed on the right is the associated negative poset $T_3^-$. Indeed, the order induced on $A$ is unchanged, while $5 \leq_-a$ for every $a \in A$. By Corollary  \ref{The associated universally derived equivalent poset}, $T_3$ and $T_3^-$ are universally derived equivalent.

The underlying undirected Hasse graph of $T_3$ is a cycle, whereas that of $T_3^-$ is the Dynkin tree $D_5$. By Proposition \ref{Reversibility of the construction}, however, the original Tamari lattice $T_3$ can be reconstructed from $T_3^-$ together with the partition $A \sqcup B$. Indeed, the unique element $5 \in B$ is the maximal element of $B$ lying below every $a \in A$ in $T_3^-$, so that $M(a)=\{5\}$ is recovered for every $ a \in A$.
\end{example}

\begin{remark}
    In Definition \ref{def admissible cuts}, condition $(c)$ does not follow from conditions $(a)$ and $(b)$. Let $B=\{u,v,w\}$, where $v < w$ and $u$ is incomparable with both $v$ and $w$, and let $A=\{a,a'\}$ with $a <a'$. Let $P=A \sqcup B$ be the poset obtained by retaining the orders on $A$ and $B$ and adding the relations $a<u$, $a<v$ and $a'<w$. Then $A$ is an ideal and $B$ is a filter, so condition $(a)$ holds. Moreover, $M(a)=\{u,v\}$, while $M(a')=\{w\}$. Condition $(b)$ holds, whereas condition $(c)$ fails for $u \in M(a)$, since $u \not\leq_B w$. 
    \\The same example shows that condition $(c)$ is needed not only to recover the bijection required in condition $(ii)$ of Section \ref{Ladkani's construction}, but also to ensure that the candidate negative relation defines a partial order: without $(c)$, transitivity can fail. Indeed, define a relation $\preceq$ on $A \sqcup B$ by retaining the orders on $A$ and $B$, with the only mixed relations given by \[b \preceq x \iff \text{ there exists } m \in M(x) \text{ such that } b \leq_B m,\] for $x \in A$ and $b \in B$. Then 
    \[u \preceq a \preceq a',\] while $u \npreceq a'$, since $M(a')=\{w\}$ and $u \nleq_Bw$. Thus $\preceq$ is not transitive and does not define a partial order.
\end{remark}
\\[-0.1 cm]
We conclude this subsection with another direct application of Theorem \ref{generalizzazione teomio}. For a finite poset $P$, let $P^{\operatorname{op}}$ denote the poset that has the same underlying set as $P$, with
\[x \leq_{P^{\operatorname{op}}}y \iff y \leq_Px.\]
Recall that a finite poset $P$ has height at most one when it contains no chain $x<y<z$ of three distinct elements. Equivalently, every element of $P$ is minimal or maximal.
A subset $C$ of a poset is called an antichain if no two distinct elements of $C$ are comparable.

For diagrams of finite-dimensional vector spaces over a field, the corresponding equivalence between bounded derived categories already follows from Ladkani \cite[Corollary 4.18]{ladkani3}. The following corollary gives the universal version directly from Theorem \ref{generalizzazione teomio}.

\begin{corollary} %[Two-level posets and opposites] 
\label{Two-level posets and opposites}
    Let $P$ be a finite poset of height at most one. Then $P$ and $P^{\operatorname{op}}$ are universally derived equivalent.

    More precisely, if $A=\operatorname{Min}(P)$, $B=P \setminus A$, then $P=A \sqcup B$ is an admissible cut and the associated negative poset is $P^-=P^{\operatorname{op}}$. 
\end{corollary}

\begin{proof}
    Since $A=\operatorname{Min}(P)$, no two distinct elements of $A$ are comparable, and hence $A$ is an antichain. Moreover, if $x \leq_Pa$ for some $a \in A$, then the minimality of $a$ implies $x=a$. Therefore $A$ is an ideal of $P$.

    We first observe that every element of $B$ is maximal in $P$. Indeed, let $b \in B$. Since $b$ is not minimal, there exists $x \in P$ such that $x<_Pb$. If $b$ were not maximal, there would also exist $y \in P$ such that $b <_Py$, giving a chain 
    \[x<_Pb<_Py\] of three distinct elements, contradicting the assumption that $P$ has height at most one. Thus every element of $B$ is maximal. It follows that $B$ is an antichain. Moreover, since every element of $B$ is maximal, $B$ is a filter of $P$. Indeed, if $b \in B$ and $b \leq_Py$, then the maximality of $b$ implies $y=b\in B$.

    Let $a \in A$. Since $B$ is an antichain, every element of $U_B(a)=\{b \in B \colon a \leq_Pb\}$ is minimal in $U_B(a)$. Therefore $M(a)=U_B(a)$. We verify the separation condition. If $m \neq m'$ belong to $M(a)$, then, since $B$ is an antichain, 
    \[[m, \bullet]_B =[\bullet, m]_B=\{m\},\] and similarly
    \[[m', \bullet]_B=[\bullet,m']_B=\{m'\}.\] Hence \[[m, \bullet]_B \cap [m', \bullet]_B= \emptyset\] and \[[\bullet, m]_B \cap [\bullet, m']_B=\emptyset.\] The compatibility condition is also automatic. Indeed, since $A$ is an antichain, if $a \leq_Aa'$, then $a=a'$. Thus, for every $m \in M(a)$, we may take $m'=m$.

    Therefore $P=A \sqcup B$ is an admissible cut.

    It remains to identify the associated negative poset. Let $a \in A$ and $b \in B$. By definition of the negative order,
    \[b \leq_-a \iff \text{ there exists } m \in M(a) \text{ such that } b \leq_Bm.\]
    Since $B$ is an antichain, the relation $b \leq_Bm$ is equivalent to $b=m$. Consequently, 
    \[b \leq_-a \iff b \in M(a).\]
    Since $M(a)=U_B(a)$, this is equivalent to $a \leq_Pb.$
    Hence \[b \leq_-a \iff a \leq_Pb.\]
    The induced orders on $A$ and $B$ are trivial, since both subsets are antichains. 
    Thus every nontrivial relation of $P$ is reversed in $P^-$, and therefore $P^-=P^{\operatorname{op}}$. 
     The universal derived equivalence between $P$ and $P^{\operatorname{op}}$ now follows from Corollary \ref{The associated universally derived equivalent poset}.
\end{proof}

\begin{example} \label{example two level posets}
Let $P$ be the poset of height at most one displayed on the left below; its opposite $P^{\operatorname{op}}$ is displayed on the right.
% https://q.uiver.app/#q=WzAsMTAsWzIsMiwiXFxidWxsZXQgYV8xIl0sWzQsMCwiXFxidWxsZXQgYl8xIl0sWzYsMiwiXFxidWxsZXQgYV8yIl0sWzQsNCwiXFxidWxsZXQgYl8yIl0sWzAsMiwiXFxidWxsZXQgYl8zIl0sWzksMiwiXFxidWxsZXQgYl8zIl0sWzExLDIsIlxcYnVsbGV0IGFfMSJdLFsxMywwLCJcXGJ1bGxldCBiXzEiXSxbMTUsMiwiXFxidWxsZXQgYV8yIl0sWzEzLDQsIlxcYnVsbGV0IGJfMiJdLFswLDFdLFsyLDFdLFsyLDNdLFswLDNdLFswLDRdLFs1LDZdLFs3LDZdLFs3LDhdLFs5LDZdLFs5LDhdXQ==
\[\begin{tikzcd}[sep=small]
	&&&& {\bullet b_1} &&&&&&&&& {\bullet b_1} && \\
	\\
	{\bullet b_3} && {\bullet a_1} &&&& {\bullet a_2} &&& {\bullet b_3} && {\bullet a_1} &&&& {\bullet a_2} \\
	\\
	&&&& {\bullet b_2} &&&&&&&&& {\bullet b_2}
	\arrow[from=1-14, to=3-12]
	\arrow[from=1-14, to=3-16]
	\arrow[from=3-3, to=1-5]
	\arrow[from=3-3, to=3-1]
	\arrow[from=3-3, to=5-5]
	\arrow[from=3-7, to=1-5]
	\arrow[from=3-7, to=5-5]
	\arrow[from=3-10, to=3-12]
	\arrow[from=5-14, to=3-12]
	\arrow[from=5-14, to=3-16]
\end{tikzcd}\]
The poset $P$ has two minimal elements $a_1,a_2$, and three maximal elements $b_1,b_2,b_3$. Notice that its underlying undirected Hasse graph contains a cycle. 

Take \[A=\{a_1,a_2\}, \,\,\,\, B=\{b_1,b_2,b_3\},\] then 
\[M(a_1)=\{b_1,b_2,b_3\}, \,\,\,\, M(a_2)=\{b_1,b_2\}.\] Thus this example uses the non-singleton case of the construction. 
By Corollary \ref{Two-level posets and opposites}, the associated negative poset is $P^{\operatorname{op}}$, obtained by reversing all the covering relations displayed above. Hence $P$ and $P^{\operatorname{op}}$ are universally derived equivalent. They are not isomorphic, since $P$ has two minimal and three maximal elements, whereas $P^{\operatorname{op}}$ has three minimal and two maximal elements.
Thus this example yields a pair of non-isomorphic universally derived equivalent posets whose common underlying undirected Hasse graph contains a cycle.
\end{example}

\subsection{The singleton case and cardinality constraints}
  
\begin{corollary} %[The singleton case] 
\label{Reconstructing a poset from a cut}
Let $P$ be a finite poset, and let $P=A \sqcup B$ be a partition of the underlying set of $P$.
The following are equivalent: \begin{enumerate}
    \item There exists an order-preserving map $f \colon A \to B$ such that $P=(A \sqcup B, \leq_f^+)$;
    \item $A$ is an ideal of $P$, $B$ is a filter of $P$, and for every $a \in A$, the set $U_B(a)$ has a minimum.
\end{enumerate}

\end{corollary} 
\begin{proof}
Assume $(ii)$. Then $M(a)$ is a singleton for every $a \in A$, say $M(a)=\{f(a)\}$. The separation condition is automatic. If $a \leq_Aa'$, then $f(a') \in U_B(a)$. Since $f(a)$ is the minimum of $U_B(a)$, we obtain $f(a) \leq_Bf(a')$. Thus $f$ is order-preserving and the compatibility condition holds. Theorem \ref{generalizzazione teomio} therefore gives $P=(A \sqcup B, \leq_f^+)$.

Conversely, assume $(i)$. Then, by the definition of $\leq_f^+$, no element of $B$ lies below an element of $A$, so $A$ is an ideal and $B$ is a filter. Moreover, for every $a \in A$, 
\[U_B(a)=\{b \in B \colon f(a) \leq_Bb\}=U_B(f(a)).\]
Hence $f(a)$ is the minimum of $U_B(a)$.

\end{proof}

We say that an admissible cut $P=A \sqcup B$ is of singleton type if $M(a)$ is a singleton for every $a \in A$.
\\[0.8 cm]
\begin{corollary} \label{cornuovo}
    Let $P=A \sqcup B$ be an admissible cut. Suppose that $P$ has a greatest element $\hat{1}$ and that $B \neq \emptyset$. Then  $M(a)$ is a singleton for every $a \in A$. Consequently, the admissible cut is of singleton type.
\end{corollary}

\begin{proof}
Since $B \neq \emptyset$ and $B$ is a filter of $P$, the greatest element $\hat{1}$ belongs to $B$. Fix $a \in A$. Then $a \leq_P \hat{1}$, and hence $\hat{1}\in U_B(a)$. Thus $U_B(a) \neq \emptyset$, so $M(a) \neq \emptyset$. 

    Suppose that $m, m' \in M(a)$. Since $ \hat{1}$ is the greatest element of $P$, we have \[m \leq_B \hat{1} \,\,\, \text{ and }\,\,\, m' \leq_B \hat{1}.\] Therefore \[\hat{1} \in [m, \bullet]_B \cap [m', \bullet]_B.\] By the separation condition of an admissible cut, this is impossible if $m \neq m'$. Hence $m=m'$, and therefore $M(a)$ is a singleton. 

    Since this holds for every $a \in A$, the admissible cut is of singleton type.
\end{proof}

\begin{remark}
    The same argument shows that $|M(a)| \leq 1$ whenever every two elements of $B$ have a common upper bound. Dually, the same conclusion holds when every two elements of $B$ have a common lower bound. In particular, this applies when $B$ is a join-semilattice or a meet-semilattice\footnote{Recall that, for two elements $x,y$ in a poset, their join $x\lor y$, when it exists, is their least upper bound, while their meet $x \land y$, when it exists, is their greatest lower bound. A poset in which every pair of elements admits a join (respectively, a meet) is called a join-semilattice (respectively, a meet-semilattice). For more details, see Garrett Birkhoff \cite[Chapter 1, page 9]{Birkhoff}.}.
\end{remark}

\begin{proposition}%[Constant cardinality on connected components] 
\label{Constant cardinality on connected components}
Let $P=A \sqcup B$ be an admissible cut. If $a, a' \in A$ belong to the same connected component of the underlying undirected Hasse graph of $A$, then \[|M(a)|=|M(a')|.\]
In particular, on each connected component $C$ of $A$, there exists an integer $r_C \geq0$ such that 
\[|M(a)|=r_C \,\,\, \text{ for every } a \in C.\]
Consequently, if $A$ is connected and $M(a_0)$ is a singleton for some $a_0 \in A$, then $M(a)$ is a singleton for every $a \in A$, and the admissible cut is of singleton type.
\end{proposition}
\begin{proof}
By Theorem \ref{generalizzazione teomio}, the sets $M(a)$ and $M(a')$ are in bijection whenever $a,a' \in A$ are comparable. Since adjacent vertices in the Hasse diagram of $A$ are comparable, the function $a \mapsto|M(a)|$ is constant on each connected component of $A$. The remaining assertions follow immediately.  
\end{proof}

\begin{remark}
    The connectedness assumption in Proposition \ref{Constant cardinality on connected components} cannot be omitted. Indeed, in Example \ref{example two level posets} the poset $A=\{a_1,a_2\}$ has two connected components, while $|M(a_1)|=3$ and $|M(a_2)|=2$. Thus the cardinality of $M(a)$ need not be constant on different connected components of $A$.
\end{remark}

\subsection{Invariance under isomorphisms}

The finite procedure described after Theorem \ref{generalizzazione teomio} may produce isomorphic outputs from different choices of data. The following result identifies a basic source of this redundancy: replacing the two input posets by isomorphic copies and transporting the construction data along the chosen isomorphisms does not change the isomorphism classes of the resulting posets. In particular, when the input posets are fixed, families lying in the same orbit under the natural action of \[\operatorname{Aut}(A) \times \operatorname{Aut}(B)\] produce isomorphic positive and negative posets. Thus, such families need not be considered separately when classifying the outcomes of the construction. 

Let $A$ and $B$ be finite posets, and let $\mathcal{Y}=(Y_a)_{a\in A}$ be a family of subsets of $B$ satisfying the hypotheses of Section \ref{Ladkani's construction}. We denote the associated positive and negative posets by \[(A \sqcup B)^+_{\mathcal{Y}} \,\,\, \text{ and }\,\,\, (A \sqcup B)^-_{\mathcal{Y}}.\] 
When it is necessary to distinguish the corresponding orders, we write $\leq^+_{\mathcal{Y}}$ and $\leq^-_{\mathcal{Y}}$.

\begin{proposition} \label{invgen}
    Let $A, B, A'$ and $B'$ be finite posets, and let $\mathcal{Y}=(Y_a)_{a \in A}$ be a family of subsets of $B$ satisfying the hypotheses of the construction in Section \ref{Ladkani's construction}. Let
    \[\alpha \colon A \to A', \,\,\,\,\,\,\, \beta \colon  B \to B'\] be isomorphisms of posets. 
    For every $a' \in A'$, define $Y'_{a'} \coloneq \beta (Y_{\alpha^{-1}(a')})$, and set $\mathcal{Y'}=(Y'_{a'})_{a' \in A'}$. 
    Then $\mathcal{Y}'$ satisfies the same hypotheses. Moreover,
    \[(A \sqcup B ) _{\mathcal{Y}}^+\cong (A' \sqcup B' )_{\mathcal{Y}'}^+\] and \[(A \sqcup B) _{\mathcal{Y}}^- \cong (A' \sqcup B')_{\mathcal{Y}'}^-.\]
\end{proposition}
\begin{proof}
    We first prove that $\mathcal{Y}'$ satisfies the same hypotheses.
    
    Let $a' \in A'$, and let $y'_1,y'_2 \in \mathcal{Y}'_{a'}$ be distinct. Set $a=\alpha^{-1}(a')$. By the definition of $Y'_{a'}$, there exist distinct elements $y_1,y_2 \in Y_a$ such that \[y'_1=\beta(y_1), \,\,\,\,\,\, y_2'=\beta(y_2).\]
    Since $\beta \colon B \to B'$ is an isomorphism of posets, it maps upper and lower cones in $B$ onto the corresponding cones in $B'$: 
    \[[\beta(y), \bullet]_{B'}=\beta([y,\bullet]_B), \,\,\,\,\,\, [\bullet, \beta(y)]_{B'}=\beta([\bullet,y]_B).\]
    Since $\mathcal{Y}$ satisfies condition (i),
    \[[y_1,\bullet]_B \cap [y_2,\bullet]_B= \emptyset \,\,\, \text{ and } \,\,\, [\bullet, y_1]_B \cap [\bullet, y_2]_B= \emptyset.\]
    Applying $\beta$, we obtain 
    \[[y'_1,\bullet]_{B'} \cap [y'_2,\bullet]_{B'}= \emptyset \,\,\, \text{ and } \,\,\, [\bullet, y'_1]_{B'} \cap [\bullet, y'_2]_{B'}= \emptyset.\]
    Thus condition (i) holds for $\mathcal{Y}'.$
    We now verify condition (ii). Let $a',a'' \in A'$ with $a' \leq_{A'} a''$. Since $\alpha \colon A \to A'$ is an isomorphism,
    \[\alpha^{-1}(a') \leq_A \alpha^{-1}(a'').\] Hence there exists a bijection 
    \[\varphi_{\alpha^{-1}(a'),\alpha^{-1}(a'')} \colon Y_{\alpha^{-1}(a')} \to Y_{\alpha^{-1}(a'')}\] such that 
    \[y \leq_B \varphi_{\alpha^{-1}(a'),\alpha^{-1}(a'')}(y)\] for every $y \in Y_{\alpha^{-1}(a')}$. Define
    \[\varphi'_{a',a''} \coloneq \beta \circ \varphi_{\alpha^{-1}(a'), \alpha^{-1}(a'')} \circ \beta^{-1}.\]
    Then $\varphi_{a',a''}' \colon Y'_{a'} \to Y'_{a''}$ is a bijection.
    
    If $z' \in Y_{a'}'$, then $z'=\beta(z)$ for some $z \in Y_{\alpha^{-1}(a')}$. Therefore,
    \[z \leq_B\varphi_{\alpha^{-1}(a'),\alpha^{-1}(a'')}(z).\] Since $\beta$ is order-preserving, $z' \leq_{B'} \varphi'_{a',a''}(z')$.
    Thus condition (ii) also holds for $\mathcal{Y}'$.

    It remains to prove the two isomorphisms. Define \[\gamma \colon A \sqcup B \to A' \sqcup B'\] by \[\gamma(a)=\alpha(a) \,\, \text{ for } a \in A, \,\,\,\,\, \gamma(b)=\beta(b) \,\, \text{ for } b \in B.\]
    Since $\alpha$ and $\beta$ are isomorphisms, $\gamma$ is a bijection and preserves the orders inside $A$ and $B$. 

    We first consider the positive mixed relations. Let $a \in A$ and $b \in B$. Then $a \leq^+_{\mathcal{Y}}b$ if and only if there exists $y \in Y_a$ such that $y \leq_B b$.

    Since $\beta \colon B \to B'$ is an isomorphism, this is equivalent to the existence of $y' \in \beta( Y_a)$ such that $y' \leq_{B'} \beta(b)$.

    Moreover, $\beta(Y_a)=Y'_{\alpha(a)}$, because \[Y'_{\alpha(a)}=\beta(Y_{\alpha^{-1}(\alpha(a))})=\beta(Y_a).\]
    Hence \[a \leq_{\mathcal{Y}}^+b \iff \alpha(a) \leq^+_{\mathcal{Y}'} \beta(b),\]

    or equivalently \[a \leq^+_{\mathcal{Y}}b \iff \gamma(a) \leq^+_{\mathcal{Y}'} \gamma (b).\]
    Therefore $\gamma$ is an isomorphism 
    \[(A \sqcup B)^+_{\mathcal{Y}} \cong (A' \sqcup B')^+_{\mathcal{Y}'}.\] 
    Similarly, for the negative mixed relations, \[b \leq_{\mathcal{Y}}^-a \iff \text{ there exists } y \in Y_a \text{ such that } b \leq_B y \iff\] \[ \text{ there exists } y' \in Y'_{\alpha(a)} \text{ such that } \beta (b) \leq_{B'} y' \iff \beta (b) \leq^-_{\mathcal{
    Y'}} \alpha(a).\] 
    Thus \[b \leq^-_{\mathcal{
    Y}}a \iff \gamma (b) \leq ^- _{\mathcal{
    Y}'} \gamma (a),\] and consequently \[(A \sqcup B)^-_{\mathcal{Y}} \cong (A' \sqcup B')^-_{\mathcal{Y}'}.\]
\end{proof}

The singleton case gives the following immediate consequence.
\begin{corollary}
    Let $A$ and $B$ be finite posets, let $f \colon A \to B$ be an order-preserving map, and let $\alpha \in \operatorname{Aut}(A)$, $\beta \in \operatorname{Aut}(B)$. 

    Define $g=\beta \circ f \circ \alpha^{-1}$. Then
    \[(A \sqcup B, \leq_f^+)\cong (A \sqcup B, \leq_g^+)\] and 
    \[(A \sqcup B, \leq_f^-) \cong (A \sqcup B, \leq_g^-).\]
\end{corollary}
\begin{proof}
    Apply Proposition \ref{invgen} with $A'=A, B'=B$, and with the singleton family $Y_a=\{f(a)\}$. Then
    \[Y'_a=\beta(Y_{\alpha^{-1}(a)})=\beta(\{f(\alpha^{-1}(a))\})=\{g(a)\}.\] The transformed family is the singleton family associated with $g$, and the isomorphisms follow from Proposition \ref{invgen}.
\end{proof}

Thus, in the singleton case, order-preserving maps which lie in the same orbit under the natural action of $\operatorname{Aut} (A)\times \operatorname{Aut}(B)$ give rise to isomorphic positive and negative posets. More generally, Proposition \ref{invgen} shows that the same invariance holds for arbitrary families satisfying the hypotheses of Section \ref{Ladkani's construction}.

\begin{remark}
    Proposition \ref{invgen} includes the case of an automorphism of the whole poset. Indeed, let $P=(A \sqcup B)_{\mathcal{Y}}^+$ and let $\sigma \in \operatorname{Aut}(P)$. Set $A'=\sigma (A)$, $B'=\sigma(B)$, and take $\alpha=\sigma \mid_A \colon A \to A'$ and $\beta=\sigma \mid_B \colon B \to B'$. Then the transformed family satisfies $Y'_{\sigma(a)}=\sigma(Y_a)$ for every $a \in A$, and Proposition \ref{invgen} gives the corresponding positive and negative isomorphisms.
\end{remark}